% 第8章后半（合并进 chapters/08_pxl.tex）
% IC Validator User Guide X-2025.06 第 8 章 PXL 后半（p211--243）
% 章节正文片段；不含 \documentclass / preamble
% 续：User-Defined Structures and Types（enum of... 跨类型说明起）

不允许跨类型（cross-type）操作。此外，在同一类型定义内，同一名称不能使用超过一次。

下例展示跨类型操作。该操作不合法：
\begin{lstlisting}[style=pxl]
v1 : t1_e;           // NONE, TOP, BOTTOM, ALL
v1 = NONE;           // Legal assignment
v1 = TOP;            // Legal assignment
v1 = RIGHT;          // Illegal assignment

v2 : t2_e;           // NONE, LEFT, RIGHT, ALL
v2 = NONE;           // Legal assignment

if (v2 == v1) // Illegal
\end{lstlisting}

\subsubsection{\cmd{list of...}}

列表（list）是对象的数组。
\begin{itemize}
  \item 所有对象属于同一类型。该类型允许为空，也可包含重复对象。
  \item 可用 \cmd{foreach()} 循环结构进行处理。
  \item 列表可用方括号（\cmd{[]}）按索引号访问。
\end{itemize}

\begin{seeAlsoBox}
关于列表的列表，见第~\pageref{sec:pxl-nested-containers}~页「嵌套容器」。
\end{seeAlsoBox}

该类型支持表~\ref{tab:pxl-list-ops} 所示操作。

\begin{longtable}{@{}>{\raggedright\arraybackslash\ttfamily}p{0.28\textwidth}
  >{\raggedright\arraybackslash}p{0.12\textwidth}
  >{\raggedright\arraybackslash}p{0.12\textwidth}
  >{\raggedright\arraybackslash}p{0.10\textwidth}
  p{0.28\textwidth}@{}}
\caption{列表操作}\label{tab:pxl-list-ops}\\
\toprule
操作 & e2 类型 & 返回 & 修改 e1 & 说明 \\
\midrule
\endfirsthead
\caption[]{列表操作（续）}\\
\toprule
操作 & e2 类型 & 返回 & 修改 e1 & 说明 \\
\midrule
\endhead
\bottomrule
\endfoot
e1.empty() & -- & boolean & 否 & 若列表不含任何元素则返回 true，否则返回 false。 \\
e1.size() & -- & integer & 否 & 返回列表中的元素个数。 \\
e1.contains(e2) & item & boolean & 否 & 若 e2 是 e1 列表的元素则返回 true。 \\
e1.front() & -- & item & 否 & 返回列表第一个元素的副本。 \\
e1.back() & -- & item & 否 & 返回列表最后一个元素的副本。 \\
e1.chop\_back() & -- & list & 否 & 返回去掉最后一个元素后的 e1 列表副本。 \\
e1.chop\_front() & -- & list & 否 & 返回去掉第一个元素后的 e1 列表副本。 \\
e1.append(e2) & item & list & 否 & 返回将 e2 追加到列表末尾后的 e1 列表副本。 \\
e1.prepend(e2) & item & list & 否 & 返回将 e2 插入到列表前端后的 e1 列表副本。 \\
e1.push\_front(e2) & item & void & 是 & 将 e2 表达式的副本添加到列表前端。 \\
e1.push\_back(e2) & item & void & 是 & 将 e2 表达式的副本添加到列表末尾。 \\
e1.pop\_front() & -- & void & 是 & 移除列表的第一个元素。 \\
e1.pop\_back() & -- & void & 是 & 移除列表的最后一个元素。 \\
e1.merge(e2) & list & list & 否 & 返回将 e2 列表拼接在末尾后的 e1 列表副本。 \\
e1.remove(e2) & item & void & 是 & 移除 e1 列表中所有等于 e2 表达式的值。 \\
\end{longtable}

下列为合法代码示例：
\begin{lstlisting}[style=pxl]
x: layer = ... ;

MyPartition: list of layer = PartitionLayer(L1);

foreach (part in MyPartition) { ... }

if (MyPartition.contains(x)) { ... }

i : integer = MyPartition.size();
\end{lstlisting}

下例展示如何拼接列表：
\begin{lstlisting}[style=pxl]
list1 : list of integer = { 1, 2, 3 };
list2 : list of integer = { 10, 20, 30 };
list3 : list of integer = list1.merge(list2);
\end{lstlisting}

此时 \cmd{list3} 的值为 \cmd{\{ 1, 2, 3, 10, 20, 30 \}}。

下面是初始化和使用字符串列表的示例：
\begin{lstlisting}[style=pxl]
// declare myList to be a list of strings, uninitialized
myList : list of string;

// declare myList to be a list of strings, initialized
myList : list of string = { "this", "is" , "a", "list" };
// myList[0] = "this"
// myList[3] = "list"
\end{lstlisting}

支持内置列表运算符函数。例如：
\begin{lstlisting}[style=pxl]
myList : list of string = { "this", "is" , "a", "list" };

s = myList.size();                 // returns the size of a list: s is "4"
a = myList.front();                // a contains a copy of L1[0]: a is "this"
myList.pop_front();                // removes the first element of the list,
                                   // the list is now { "is", "a", "list" }
myList.pop_back();                 // removes the last element of the list,
                                   // the list is now { "is", "a" }
myList.push_front("xyz");          // adds a new element containing "xyz"
                                   // to the front of the list.
                                   // the list is now { "xyz", "is", "a" }
myList.push_back("bar");           // adds a new element containing "bar"
                                   // to the back of the list.
                                   // the list is now { "xyz", "is", "a", "bar" }
\end{lstlisting}

\subsubsection{\cmd{hash of...}}

定义从一种数据类型到另一种数据类型的映射。源类型称为键类型（key type）；目标类型称为值类型（value type）。也就是说，hash 定义了键/值对数组。

该类型要求键类型为 \cmd{string}、\cmd{integer} 或 \cmd{enum}，或由这三种类型之一派生的类型。不能是 \cmd{double}。值类型可以是任何先前已定义的类型。

下例创建将整数转换为字符串的 hash。此时键类型为 \cmd{integer}，值类型为 \cmd{string}。该类型使用 \cmd{=>} 运算符初始化。
\begin{lstlisting}[style=pxl]
t1_h : newtype hash of integer to string;
t2_h : newtype hash of string to integer;

h1 : t1_h = {
   1 => "a",
   2 => "b",
   26 => "z"
};

h2 : t2_h = {
   "xyzzy" => 1,
   "plugh" => 10,
   "help" => 100
};
\end{lstlisting}

下面是使用派生类型的示例：
\begin{lstlisting}[style=pxl]
t1 : newtype integer;
h1_h : newtype hash of t1 to string;
x : h1_h = { 1 => "a" } ;
\end{lstlisting}

\cmd{hash of...} 类型可使用下标按元素赋值：
\begin{lstlisting}[style=pxl]
x = h1[203];
h1[30] = "ad";
\end{lstlisting}

\cmd{hash of...} 类型支持表~\ref{tab:pxl-hash-ops} 所示成员操作。

\begin{longtable}{@{}>{\raggedright\arraybackslash\ttfamily}p{0.28\textwidth}
  >{\raggedright\arraybackslash}p{0.14\textwidth}
  >{\raggedright\arraybackslash}p{0.12\textwidth}
  >{\raggedright\arraybackslash}p{0.10\textwidth}
  p{0.26\textwidth}@{}}
\caption{Hash 操作}\label{tab:pxl-hash-ops}\\
\toprule
操作 & e2 类型 & 返回 & 修改 e1 & 说明 \\
\midrule
\endfirsthead
\caption[]{Hash 操作（续）}\\
\toprule
操作 & e2 类型 & 返回 & 修改 e1 & 说明 \\
\midrule
\endhead
\bottomrule
\endfoot
e1.contains\_key(e2) & hash 的键类型 & boolean & 否 & 若 e2 是 e1 哈希表的键则返回 true。效果上与 \cmd{e1.keys().contains(e2)} 相同，但运行时性能可能更好。 \\
e1.contains\_value(e2) & hash 的值类型 & boolean & 否 & 若 e2 是 e1 哈希表中某项的值则返回 true。效果上与 \cmd{e1.values().contains(e2)} 相同，但运行时性能可能更好。 \\
e1.empty() & -- & boolean & 否 & 若 hash 不含任何元素则返回 true，否则返回 false。 \\
e1.keys() & -- & list of 键类型 & 否 & 返回所有键值的列表；该列表不含重复值。 \\
e1.size() & -- & integer & 否 & 返回 hash 中的元素个数。 \\
e1.values() & -- & list of 值类型 & 否 & 返回所有值的列表；根据 hash 内容，该列表可含重复值。 \\
e1.merge(e2) & 与 e1 匹配的 hash 类型 & hash & 否 & 返回将 e2 的键与值拼接到末尾后的 e1 hash 副本；若某键在 e1 与 e2 中均存在，则取其在 e2 中的值。 \\
e1.remove(e2) & item & void & 是 & 从 e1 hash 中移除 e2 键。若 e2 键不存在则无效果。 \\
\end{longtable}

再举一例：
\begin{lstlisting}[style=pxl]
alpha_h : newtype hash of string to integer;
alpha : alpha_h = {       // declare and initialize key => value pairs
        "a" => 1,
        "b" => 2,
        "c" => 3
        };
y = alpha["b"];          // The value of y is "2"
keyList = alpha.keys(); // A list of the keys in alpha
\end{lstlisting}

\subsubsection{\cmd{struct of...}}

定义由用户提供名称标识的独立类型。必须至少有一个成员，且成员名必须唯一。各成员可以属于不同的类型。

下面是一个用几何坐标定义点的结构示例：
\begin{lstlisting}[style=pxl]
point: newtype struct of { x: integer; y: integer; };
\end{lstlisting}

\cmd{struct of...} 类型可包含初始值，用作结构创建时的默认值。
\begin{lstlisting}[style=pxl]
t3_s : newtype struct of {
   x: integer;
   y: integer = 76;   // initializer
   z: string = "hi"; // initializer
};
t4_s : newtype struct of {
   a: integer = 106; // initializer
   b: string = "boo"; // initializer
};
\end{lstlisting}

\cmd{struct of...} 类型可从 C 风格初始化器赋值，也可从名--值对列表赋值。下例展示使用上例所定义类型的赋值。
\begin{lstlisting}[style=pxl]
s3 : t3_s = { 10, z = "bye" };            // note that 'y' defaults
                                          // to the value 76

// the defaults cause s4 to have a=106, b="boo"
s4 : t4_s = { };
\end{lstlisting}

再举一例：
\begin{lstlisting}[style=pxl]
shoebox : newtype struct of {
      style : string = "Tennis Shoe";
      shoe_size : double = 10.5;
   };

box : shoebox;
shoe = box.style;          // shoe is "Tennis shoe", the default
box.style = "Dress Shoe"; // Reassign the member box.style
// create a new shoebox with a different kind of shoe
box2 : shoebox = {"Boat Shoe", 9};
\end{lstlisting}

\subsubsection{\cmd{function}}

允许将函数传递给另一个用户函数。被传递的函数可像全局函数一样使用。例如：
\begin{lstlisting}[style=pxl]
// input file
#include <icv.rh>
xyzzy : function (void) returning void { note("abcd"); }
plugh : function ( f : function(void) returning void ) returning void {
f(); }
plugh(xyzzy); // summary file shows "abcd" in output when plugh runs
\end{lstlisting}

\subsubsection{嵌套类型}
\subsubsection*{Nested Types}

类型可以组合形成嵌套类型，例如：
\begin{itemize}
  \item \cmd{list of struct}
  \item \cmd{list of enum}
\end{itemize}

例如，使用上文 \cmd{struct of...} 中 \cmd{shoebox} 结构定义：
\begin{lstlisting}[style=pxl]
box : list of shoebox = {
         { "Walking", 12 },
         { "High Heel", 7 }
      };

A = box[0].style;                  // A is "Walking"
B = box[1].shoe_size;              // B is 7
\end{lstlisting}

\subsection{类型分析}
\subsubsection*{Type Analysis}

本节说明 PXL 在特殊情况下如何处理数据类型，例如当参数传入一种数据类型但期望另一种数据类型时。

当 PXL 遇到一种数据类型而期望另一种时，只要合适就会执行自动类型转换。用于自动类型转换的规则在下列各小节中讨论。

\subsubsection{标准算术转换}
\subsubsection*{Standard Arithmetic Conversions}

当表达式所要求的算术类型与实际找到的算术类型不匹配时，按表~\ref{tab:pxl-arith-conv} 进行转换。

\begin{longtable}{@{}p{0.26\textwidth}
  >{\centering\arraybackslash}p{0.14\textwidth}
  >{\centering\arraybackslash}p{0.16\textwidth}
  >{\centering\arraybackslash}p{0.14\textwidth}
  >{\centering\arraybackslash}p{0.16\textwidth}@{}}
\caption{算术转换}\label{tab:pxl-arith-conv}\\
\toprule
期望 \ldots & 允许 integer\textsuperscript{a} & 允许 constraint of integer\textsuperscript{b} & 允许 double & 允许 constraint of double \\
\midrule
\endfirsthead
\caption[]{算术转换（续）}\\
\toprule
期望 \ldots & 允许 integer & 允许 constraint of integer & 允许 double & 允许 constraint of double \\
\midrule
\endhead
\bottomrule
\endfoot
integer & 是 & 否 & 否 & 否 \\
constraint of integer & 是 & 是 & 否 & 否 \\
double & 是 & 否 & 是 & 否 \\
constraint of double & 是 & 是 & 是 & 是 \\
\end{longtable}

{\footnotesize
\textsuperscript{a}「是」表示允许该转换。\\
\textsuperscript{b}「否」表示不允许该转换。该操作将导致错误。
}

这些转换仅在表达式类型要求时才应用。特别地，除非上下文要求，否则 \cmd{integer} 类型的表达式不会转换为 \cmd{double}。

将 \cmd{integer} 或 \cmd{double} 转换为 constraint 时，所得 constraint 看起来如同使用 \cmd{==} 一元约束（\cmd{CONSTRAINT\_EQ}）书写一样。
\begin{lstlisting}[style=pxl]
f: function ( c: constraint ) returning void { ... }

f( (1.0, 26.0) );                     // constraint is ok
f( 2.0 );                             // single value is ok also
f( 42 );                              // single integer is accepted
\end{lstlisting}

\subsubsection{比较 double}
\subsubsection*{Comparing Doubles}

由于类型 \cmd{double} 是实数的近似，偶尔会出现累积误差，使表面上等价的运算产生不同结果。PXL 提供特殊函数，用于在考虑此类微小误差的情况下比较 double。这些特殊函数见 IC Validator Reference Manual 中「Utility Functions」一章 Math Functions 节的 Equivalence Comparison Functions 与 Double Comparison Functions。

\begin{noteBox}
只要涉及与 double 相关类型之一的直接比较（\cmd{==} 或 \cmd{!=}），就会报告警告。
\end{noteBox}

\subsubsection{类型提升}
\subsubsection*{Type Promotion}

当期望列表，而找到的表达式类型与该列表所含元素类型相同时，该表达式会从单个项提升为仅含该项的列表。

\subsubsection{参数传递}
\subsubsection*{Parameter Passing}

带有 \cmd{output} 或 \cmd{input-output} 限定的函数参数，其实参必须匹配。前述规则不适用。例如，当函数期望标记为 \cmd{out} 的 constraint 参数时，提供 \cmd{integer} 类型变量名不合法，因为所需转换不合法。

\subsubsection{字符串拼接}
\subsubsection*{String Concatenation}

表~\ref{tab:pxl-str-concat} 所示类型的任何表达式，均可使用标准加法运算符（\cmd{+}）追加到字符串表达式。

\begin{longtable}{@{}>{\raggedright\arraybackslash\ttfamily}p{0.32\textwidth} p{0.60\textwidth}@{}}
\caption{字符串拼接}\label{tab:pxl-str-concat}\\
\toprule
类型 & 追加到结果字符串的内容 \\
\midrule
\endfirsthead
\caption[]{字符串拼接（续）}\\
\toprule
类型 & 追加到结果字符串的内容 \\
\midrule
\endhead
\bottomrule
\endfoot
integer & 整数表达式的值 \\
double & double 表达式的值 \\
handle & handle 的文本表示；每个不同对象唯一 \\
string & 字符串内容 \\
boolean & \cmd{true} 或 \cmd{false} \\
constraint of integer & 等价的 constraint 字面量表达式 \\
constraint of double & 等价的 constraint 字面量表达式 \\
user-defined enum & 该表达式所表示的 enum 字面量名称 \\
\end{longtable}

使用字符串时请记住：
\begin{itemize}
  \item 当加法运算符（\cmd{+}）左侧为字符串时，结果为字符串。
  \item 对字符串值表达式，加法运算符（\cmd{+}）不对称：\cmd{a + b != b + a}。
  \item 后续追加到字符串的各项之间不添加空格，因此可能看到意外结果，如下例：
\begin{lstlisting}[style=pxl]
s : string;
s = "hello " + 10 + " " + (1 == 0);
// s now contains the string "hello 10 false"

c : constraint = ( 10, 40 ];
s = "" + c;
// s now contains the string "(10.,40.]", because the
// type of "c" is constraint of double

mye : newtype enum of { A, B, C, D, X };
t : mye = X;
s = "" + t + A + B;
// s now contains the string "XAB"

s = "" + 11 + 3.14 + (1 == 1);
// s now contains the string "113.14true"

s = "" + (11 + 3.14) + (1 == 1);
// s now contains the string "14.14true"
\end{lstlisting}
\end{itemize}

\subsubsection{默认值注入}
\subsubsection*{Default Injection}

处理函数调用时，未给出值的参数若有默认值，则采用函数原型中的默认值。既无默认值也未指定值的参数会导致错误。

类似地，处理字面量结构时，未给出值的成员若有默认值，则采用 \cmd{newtype} 结构定义中的默认值。既无默认值也未指定值的成员会导致错误。

\subsection{嵌套容器}
\subsubsection*{Nested Containers}
\label{sec:pxl-nested-containers}

列表、hash 与结构也称为容器（container）。容器可以多种不同组合嵌套。嵌套容器的一个例子是结构的列表。

下面是 hash 的 hash 示例。runset 代码为（hash of a hash）：
\begin{lstlisting}[style=pxl]
hash1_h : newtype hash of string to integer;
hash2_h : newtype hash of string to hash1_h;

multi_hash : hash2_h = {};

multi_hash["A"] = {};                 // Initialize hash of "A" to empty
multi_hash["A"]["ONE"] = 1;
multi_hash["A"]["TWO"] = 2;
multi_hash["B"] = {
   "THREE" => 3,
   "FOUR" => 4
};

foreach(key1 in multi_hash.keys() ){
   foreach(key2 in multi_hash[key1].keys()){
      note( "multi_hash[" + key1 + "][" + key2 + "] = "
            + multi_hash[key1][key2] + "\n");
   }
}
\end{lstlisting}

输出为：
\begin{lstlisting}
runset.rs:123 multi_hash[A][ONE] = 1
runset.rs:123 multi_hash[A][TWO] = 2
runset.rs:123 multi_hash[B][FOUR] = 4
runset.rs:123 multi_hash[B][THREE] = 3
\end{lstlisting}

下面是列表的列表示例。runset 代码为（list of a list）：
\begin{lstlisting}[style=pxl]
list1_l : newtype list of integer;
list2_l : newtype list of list1_l;

list2 : list2_l = {
   { 0, 1, 2, 3 },
   { 10, 11, 12, 13 }
   };

list2.push_back({ 20, 21, 22, 23 });
list2[1].push_back( 14 );

for( j=0 to (list2.size() - 1) ){
   for( i=0 to (list2[j].size() - 1) ){
      note( "list2[" + j + "][" + i + "] = " + list2[j][i] + "\n");
   }
}
\end{lstlisting}

输出为：
\begin{lstlisting}
runset.rs:143 list2[0][0] = 0
runset.rs:143 list2[0][1] = 1
runset.rs:143 list2[0][2] = 2
runset.rs:143 list2[0][3] = 3
runset.rs:143 list2[1][0] = 10
runset.rs:143 list2[1][1] = 11
runset.rs:143 list2[1][2] = 12
runset.rs:143 list2[1][3] = 13
runset.rs:143 list2[1][4] = 14
runset.rs:143 list2[2][0] = 20
runset.rs:143 list2[2][1] = 21
runset.rs:143 list2[2][2] = 22
runset.rs:143 list2[2][3] = 23
\end{lstlisting}

再举一个列表的列表示例。runset 代码为：
\begin{lstlisting}[style=pxl]
list3 : list of list of integer = {};

list3.push_back( { 0, 1, 2, 3 } );
list3.push_back( { 4, 5, 6 } );
list3.push_back( { 7, 8 } );
list3.push_back( { 9 } );

for( j=0 to (list3.size() - 1) ){
   for( i=0 to (list3[j].size() - 1) ){
     note( "list3[" + j + "][" + i + "] = " + list3[j][i] + "\n");
   }
}
\end{lstlisting}

输出为：
\begin{lstlisting}
runset.rs:159 list3[0][0] = 0
runset.rs:159 list3[0][1] = 1
runset.rs:159 list3[0][2] = 2
runset.rs:159 list3[0][3] = 3
runset.rs:159 list3[1][0] = 4
runset.rs:159 list3[1][1] = 5
runset.rs:159 list3[1][2] = 6
runset.rs:159 list3[2][0] = 7
runset.rs:159 list3[2][1] = 8
runset.rs:159 list3[3][0] = 9
\end{lstlisting}

% ============================================================
\section{流程控制}
\subsection*{Flow Control}

流程可用下列方式控制：
\begin{itemize}
  \item 条件语句（Conditional Statements）
  \item 循环（Loops）
  \item 执行终止（Termination of Execution）
\end{itemize}

\subsection{条件语句}
\subsubsection*{Conditional Statements}

PXL 中的条件语句是一类构造，允许你根据指定条件求值为 true 还是 false 来执行不同动作。

\cmd{if} 语句用于仅在预定义条件满足时执行某动作。

\cmd{if} 语句至少与另外两个组成部分一起使用：
\begin{itemize}
  \item 一个可求值为布尔值 true 或 false 的表达式。
  \item 一个在求值表达式返回 true 时执行的组成部分。
\end{itemize}

可选地，在必需组成部分之外，还可包含下列组成部分以构建更复杂的条件语句：
\begin{itemize}
  \item \cmd{else} 构造，后跟在第一个表达式求值为 false 时执行的组成部分。
  \item \cmd{elseif} 构造，表示在父 \cmd{if} 条件语句内嵌套了循环式 \cmd{if} 条件或一个或多个 \cmd{if} 条件。PXL 也支持 \cmd{elif} 构造。
\end{itemize}

\cmd{if} 语句使用下列语法：
\begin{lstlisting}[style=pxl]
if (e1) s1;
if (e1) s1 elif (e2) s2;
if (e1) s1 elseif (e2) s2;
if (e1) s1 else s2;
\end{lstlisting}

\cmd{if} 语句按如下方式求值：
\begin{itemize}
  \item 始终先求值第一个布尔表达式（\cmd{e1}）。
  \item 若 \cmd{e1} 求值为 true，则处理第一个组成部分（\cmd{s1}）。
  \item 若遇到 \cmd{else} 语句，则执行其后的语句（\cmd{s2}）。
  \item 若遇到 \cmd{elif} 或 \cmd{elseif} 语句，则执行其后的语句。
\end{itemize}

\subsection{循环}
\subsubsection*{Loops}

与其他编程语言一样，可在 PXL 中使用循环将某动作重复若干次。

PXL 提供下列方法创建循环构造：
\begin{itemize}
  \item \cmd{for} 循环
  \item \cmd{foreach} 循环
  \item \cmd{while} 循环
\end{itemize}

\subsubsection{\cmd{for} 循环}
\subsubsection*{for Loop}

\cmd{for} 循环用于将特定动作重复指定次数。它至少包含一个整数（\cmd{e1}）以及必须重复的动作（\cmd{s1}）。也可包含附加整数（\cmd{e2}、\cmd{e3}）。

语法为：
\begin{lstlisting}[style=pxl]
for (var = e1 to e2) s1;         //s1 is a statement; it can be compound
for (var = e1 to e2 step e3) s1;
\end{lstlisting}

\begin{itemize}
  \item 要求 \cmd{e1}、\cmd{e2} 与 \cmd{e3} 为整数值表达式。
  \item 将 \cmd{var} 定义为仅在语句 \cmd{s1} 内可见的整数变量。
  \item 在循环开始时对 \cmd{e1}、\cmd{e2} 与 \cmd{e3} 各求值一次；随后对这些表达式取值的赋值不影响循环。
\end{itemize}

例如：
\begin{lstlisting}[style=pxl]
for(i=0 to 10 step 2){              // Increment by 2: 0, 2, 4, 6, . . .
   if( i < 6) {
      a[i] = b[i] + c[i];
   } elif( ( i >=6 ) && ( i < 10)){
      a[i] = b[i] + c[i] + 1;
   } else {
      a[i] = b[i] - c[i];
   }
};
\end{lstlisting}

\subsubsection{\cmd{foreach} 循环}
\subsubsection*{foreach Loop}

\cmd{foreach} 循环类似于 \cmd{for} 循环，但用于针对容器（如列表或 enum）的元素执行一组动作。它不使用变量计数器来决定动作的迭代次数。

它包含类型为列表的 \cmd{e1}，以及对 \cmd{e1} 中每个元素求值的组成部分 \cmd{s1}。

语法为：
\begin{lstlisting}[style=pxl]
foreach (var in e1) s1;         //e1 is of the type "list of"
\end{lstlisting}

\begin{itemize}
  \item 对 \cmd{e1} 的每个元素求值一次 \cmd{s1}，同时将 \cmd{var} 从列表复制过来。
  \item 若输入列表 \cmd{e1} 为空，则跳过 \cmd{s1}。
  \item \cmd{var} 是类型正确的常量变量，仅在语句 \cmd{s1} 内可见。
\end{itemize}

\subsubsection{\cmd{while} 循环}
\subsubsection*{while Loop}

\cmd{while} 循环用于重复特定动作直到满足某条件。它包含布尔表达式（\cmd{e1}）以及必须重复的动作（\cmd{s1}）。\cmd{while} 循环始终至少对布尔表达式 \cmd{e1} 求值一次。

语法为：
\begin{lstlisting}[style=pxl]
while (e1) s1;       // e1 is evaluated as a Boolean;
                     // loop continues until e1 is false
\end{lstlisting}

\cmd{while} 语句按如下方式求值：
\begin{itemize}
  \item 每次 \cmd{e1} 求值为 true 时，执行 \cmd{s1}。
  \item 当 \cmd{e1} 求值为 false 时，该语句结束。
\end{itemize}

\subsection{执行终止}
\subsubsection*{Termination of Execution}

\cmd{error("message")} 是诊断函数（Diagnostics Function）。它写出所提供的消息，然后终止执行。例如：
\begin{lstlisting}[style=pxl]
#ifdef ABC
#ifdef DEF
error("ABC and DEF are mutually exclusive. Please do not set both.");
#endif
#endif
\end{lstlisting}

% ============================================================
\section{函数}
\subsection*{Functions}

PXL 中的函数按用途分类。函数类型见表~\ref{tab:pxl-func-types}。

\begin{longtable}{@{}>{\raggedright\arraybackslash}p{0.22\textwidth} p{0.70\textwidth}@{}}
\caption{函数类型}\label{tab:pxl-func-types}\\
\toprule
函数类型 & 说明 \\
\midrule
\endfirsthead
\caption[]{函数类型（续）}\\
\toprule
函数类型 & 说明 \\
\midrule
\endhead
\bottomrule
\endfoot
Runset function & IC Validator 发布的标准函数。 \\
Remote function & 由用户编写、可被 runset 函数调用的函数。IC Validator 工具也提供默认 remote 函数。 \\
User function & 由用户编写的函数。 \\
Utility function & 由 IC Validator 工具提供、用于在 remote 函数中访问原始数据的函数。 \\
\end{longtable}

\subsection{函数定义}
\subsubsection*{Function Definitions}

函数定义为一组与名称关联的操作。函数定义标识所有参数以及返回值的名称、类型与默认值。

你可以创建自己的函数，并与所提供的函数一起在 runset 中使用。

如下例所示，所有函数定义必须位于顶层。
\begin{lstlisting}[style=pxl]
// a simple function definition

myinversion: function ( arg: double ) returning rv : double
{
   rv = 1.0 / arg;
}

x = myinversion(2.0);         // simple function call
\end{lstlisting}

函数定义为所有处理提供单独作用域。作用域由书写位置的上下文决定（静态作用域，static scoping）。

\begin{seeAlsoBox}
关于作用域，见「变量作用域」（原书第~197~页）。
\end{seeAlsoBox}

在函数定义内部适用下列作用域规则：
\begin{enumerate}
  \item 若在当前作用域找不到声明，则向上搜索更高作用域，直到搜遍整个函数定义，在匹配处停止。
  \item 然后，如有必要，使用在函数定义之前声明且同名的全局变量。
  \item 最后，若该变量正被赋值为函数调用的结果，或处于函数调用的输出参数位置，则在当前语句的作用域中隐式声明它。但是，若该变量处于函数调用的输入或输入--输出参数位置，即使它正被赋值为函数调用的结果，也不会在当前语句的作用域中隐式声明。
\end{enumerate}

在函数定义外部适用下列作用域规则：
\begin{enumerate}
  \item 若在当前作用域找不到声明，则向上搜索更高作用域，直到搜遍整个 runset，在匹配处停止。
  \item 然后，如有必要，使用先前声明且同名的全局变量。
  \item 最后，若该变量正被赋值为函数调用的结果，或处于函数调用的输出参数位置，则在当前语句的作用域中隐式声明它。但是，若该变量处于函数调用的输入或输入--输出参数位置，即使它正被赋值为函数调用的结果，也不会在当前语句的作用域中隐式声明。
\end{enumerate}

\subsection{函数原型}
\subsubsection*{Function Prototypes}

函数原型定义函数的参数以及可选参数的默认值。函数原型标识按引用调用（call-by-reference）的参数。函数原型不定义函数的操作。

下面是函数原型示例：
\begin{lstlisting}[style=pxl]
my_xor : function (
   a : polygon_layer,
   b : polygon_layer
) returning xorOut: polygon_layer;
\end{lstlisting}

下面是函数声明示例：
\begin{lstlisting}[style=pxl]
my_xor : function (
   a : polygon_layer,
   b : polygon_layer
) returning xorOut: polygon_layer {
      tmp1 = not(a,b);
      tmp2 = not(b,a);
      xorOut = or(tmp1, tmp2 );
   };
\end{lstlisting}

下例展示调用前一例中定义的 \cmd{my\_xor()} 函数：
\begin{lstlisting}[style=pxl]
a1 : polygon_layer;
b1 : polygon_layer;
xorA1B1 = my_xor(a1, b1);
\end{lstlisting}

\subsection{函数调用}
\subsubsection*{Function Calls}

函数调用定义函数如何运作。函数调用由函数名后跟通过实参说明如何及在何处执行该函数的指令组成。

向函数调用的实参传递值可用两种不同风格：
\begin{itemize}
  \item 仅按实参值（按值传递，pass-by-value）。值的顺序必须与函数原型中参数顺序匹配。使用此风格时，若你关心某个可选参数，则必须为其前面的所有可选参数指定值，因为值按顺序匹配到参数。
  \item 按 \cmd{argument\_name = argument\_value} 格式用名称（按名传递，pass-by-name）。参数顺序无关。使用此风格时，无需为其前面的所有可选参数指定值，因为值按名称匹配到参数。
\end{itemize}

两种风格可在同一次函数调用中使用。但是，在函数调用中，所有按值传递风格指定的实参值必须出现在所有按名传递风格指定的实参值之前。一旦在函数调用中使用了按名传递风格，其后的所有实参调用必须使用相同风格。

例如，\cmd{interacting()} 函数的语法为：
\begin{lstlisting}[style=pxl]
interacting : binary published function(
   layer1 : polygon_layer,
   layer2 : polygon_layer,
   count : count_constraint_t = >0,
   include_touch : more_include_touch_e = EDGE,
   processing_mode : processing_mode_e = HIERARCHICAL
) returning result : polygon_layer;
\end{lstlisting}

下面是以按名传递风格调用 \cmd{interacting()} 的示例：
\begin{lstlisting}[style=pxl]
y = interacting( layer1 = a, layer2 = b, include_touch = ALL );
   // The count argument is skipped, and the value is the default.
   // The value of the processing_mode argument is the default.
\end{lstlisting}

下面是以按值传递风格调用 \cmd{interacting()} 的示例：
\begin{lstlisting}[style=pxl]
y = interacting( a, b, 2 );
   // The value of count is 2.
   // The values of the include_touch and processing_mode arguments
   // are the defaults.
\end{lstlisting}

下面是两种混合风格调用 \cmd{interacting()} 的示例：
\begin{lstlisting}[style=pxl]
y = interacting( a, b, include_touch = ALL );
   // The count argument is skipped, and the value is the default.
   // The value of the processing_mode argument is the default.

y = interacting( a, b, count = 2);
   // The values of the include_touch and processing_mode arguments
   // are the defaults.
\end{lstlisting}

函数调用有两种风格：
\begin{itemize}
  \item 标准函数调用（Standard Function Call）
\begin{lstlisting}[style=pxl]
y = and(a, b);
y = and(a, b, CELL_LEVEL);
y = or(not(a,b),not(b,a));
\end{lstlisting}
  \item 运算符风格函数调用（Operator-Style Function Call）
\begin{lstlisting}[style=pxl]
y = a and b;
y = a and[CELL_LEVEL] b;
y = (a not b) or (b not a);
\end{lstlisting}
\end{itemize}

\subsubsection{标准函数调用}
\subsubsection*{Standard Function Call}

对任何函数始终允许标准函数调用。

标准函数调用由函数名后跟括号括起、逗号分隔的至少一个表达式列表，再跟至少一个名--值对列表组成。

下例中函数有三个参数：
\begin{lstlisting}[style=pxl]
g: function ( x : integer, y : integer, z : integer)
   returning void;
\end{lstlisting}

下列四个函数调用中，有一个不合法：
\begin{lstlisting}[style=pxl]
g(x=3, y=10, z=10);          // allowed

g(z=3, y=10, x=4);           // allowed

g(10, 20, 30);               // allowed; x=10, y=20, z=30

g(z=3, 10, 20);              // ILLEGAL - expressions must not
                             // follow name-value pairs
\end{lstlisting}

若任一参数存在默认值，则不必指定这些参数。下例中函数有三个参数，其中两个有默认值。
\begin{lstlisting}[style=pxl]
f: function
    ( x : integer,
      y : integer = 27,
      z : integer = 42
    ) returning void { ... }
...
\end{lstlisting}

下列五个函数调用中，有两个不合法：
\begin{lstlisting}[style=pxl]
f(10, y=2);          // allowed, z is 42

f(10);               // allowed, y is 27, z is 42

f(z=3, x=4);         // allowed, y is 27

f();                 // ILLEGAL - x requires a value

f(z=20);             // ILLEGAL - x requires a value
\end{lstlisting}

\subsubsection{运算符风格函数调用}
\subsubsection*{Operator-Style Function Call}

运算符风格函数调用要求二元函数在函数名左右两侧各有实参。其他实参须以名--值对形式写在紧随函数名之后的方括号中。

仅对以 \cmd{binary} 限定符声明或定义的函数允许运算符风格函数调用。使用运算符风格函数调用有助于防止在其他上下文（如变量名与成员名）中使用该函数名。

下列语法定义如何使用 \cmd{binary} 限定符调用函数：
\begin{lstlisting}[style=pxl]
length: binary function (l1: layer, c1: constraint of double) returning
result: layer {...}
...
LONGMET = MET1 length > 5;
\end{lstlisting}

通过仔细的接口定义，\cmd{binary} 可写出易读的程序。例如：
\begin{lstlisting}[style=pxl]
// A simple declaration that takes two arguments.
// The LHS is the layer to be selected from; the RHS is the
// constraint defining the criteria area.
...
// ... as is a variable of appropriate type.
c : constraint = (27.0, 42.5);
x : lyr1 area c;
\end{lstlisting}

\subsubsection{实参绑定}
\subsubsection*{Argument Bindings}

实参绑定决定调用点的哪些实参绑定到函数原型中的哪些参数。
\begin{itemize}
  \item 根据函数名以及上下文所要求的类型选择候选原型。若无法确定类型，则仅使用函数名。进行运算符风格函数调用时，外部实参绑定到原型中的参数。对二元函数，左侧实参绑定到第一个参数，右侧实参绑定到第二个参数。
  \item 每个名--值实参绑定到剩余的同名参数。
  \item 任何没有名称的实参从左到右绑定到剩余的无默认值参数。
  \item 若所有剩余未绑定参数都有默认值，则使用该函数。
  \item 若实参绑定未能恰好匹配一个函数原型，则为错误。
\end{itemize}

例如，给定如下定义的函数 \cmd{MyFunc}：
\begin{lstlisting}[style=pxl]
MyFunc: binary function (
   fromLayer: lhs layer,
   touching: rhs layer,
   corners: boolean = true,
   whole_polygon: boolean
) returning e: edge_layer;
\end{lstlisting}

下列函数调用全部等价：
\begin{lstlisting}[style=pxl]
T1 = MyFunc (metalLayer, selectLayer, true);
T2 = MyFunc (metalLayer, selectLayer, whole_polygon=true);
T3 = metalLayer MyFunc[whole_polygon=true] selectLayer;
\end{lstlisting}

下列代码因 \cmd{whole\_polygon} 实参没有默认值而产生编译时错误：
\begin{lstlisting}[style=pxl]
/* illegal function call */
Tbroken = metalLayer MyFunc selectLayer;
\end{lstlisting}

\subsubsection{默认值}
\subsubsection*{Defaults}

默认值在函数原型中指定为选项变量的右侧赋值。例如，\cmd{and()} 函数的语法为：
\begin{lstlisting}[style=pxl]
and : function(
   layer1 : polygon_layer,
   layer2 : polygon_layer,
   processing_mode : processing_mode_e = HIERARCHICAL,
   remove_hierarchical_overlap : boolean = true
) returning and_result : polygon_layer;
\end{lstlisting}

下列对 \cmd{and()} 的调用对 \cmd{processing\_mode} 与 \cmd{remove\_hierarchical\_overlap} 实参使用默认值：
\begin{lstlisting}[style=pxl]
c = and( a, b );
\end{lstlisting}

下列对 \cmd{and()} 的调用对 \cmd{processing\_mode} 使用可选设置，对 \cmd{remove\_hierarchical\_overlap} 使用默认值：
\begin{lstlisting}[style=pxl]
c = and( a, b, CELL_LEVEL );
\end{lstlisting}

\paragraph{重置默认值}
\paragraph*{Resetting Defaults}

你可以覆盖任何函数实参的默认值。任何实参都可更改其默认值；可以是数值、布尔值或枚举量。

本例中设置了 \cmd{abc} 函数实参 \cmd{x} 的默认值：
\begin{lstlisting}[style=pxl]
abc : literal function ( x : double = 123) returning z : double;
A = abc();   // == abc(123)
\end{lstlisting}

\begin{noteBox}
警告：重置默认值时，你是在重新定义该函数，必须列出其全部实参。
\end{noteBox}

\subsection{匿名函数}
\subsubsection*{Anonymous Functions}

若干 IC Validator 命令允许将用户函数作为实参传入。匿名函数为简单行为或组织更复杂行为提供了一种机制。可在期望函数名的地方使用匿名函数。编译器生成内部名称，并在紧临其调用之前的代码位置声明该函数。

传统上，可使用全局变量，并按需为不同调用设置不同值。例如：
\begin{lstlisting}[style=pxl]
x : integer;
y : integer;
calc_nmos_props : function (void) returning void { ... }

x = 0;
y = 5;
nmos(..., calc_nmos_props, ...);

x = 1;
y = 4;
nmos(..., calc_nmos_props, ...);
\end{lstlisting}

然而，匿名函数可使此类代码更紧凑、组织更好。下例中，在每次 NMOS 调用之前声明一个隐藏函数，随后使用它。
\begin{lstlisting}[style=pxl]
calc_nmos_props : function (x : integer, y : integer) returning void
 { ... }

nmos(..., function { calc_nmos_props(0, 5); }, ...);
nmos(..., function { calc_nmos_props(1, 4); }, ...);
\end{lstlisting}

\subsection{函数重载}
\subsubsection*{Function Overloading}

在某些情况下，PXL 允许存在多个同名函数。若 PXL 能根据类型与参数名确定要调用哪个函数，则多个函数可以同名。若 PXL 无法确定要调用的函数，则导致错误。

若一个返回值而另一个返回 \cmd{void}，则可以有两个同名的不同函数。

当函数出现在语句级时，假定为返回 \cmd{void} 的函数；当出现在表达式中时，假定为返回值的函数：
\begin{lstlisting}[style=pxl]
MyFunction(a,b,c);             // use the void-returning function
                               // "MyFunction"
x = MyFunction(a,b,c);         // use the value-returning function
                               // "MyFunction"
\end{lstlisting}

\begin{noteBox}
PXL 仅支持基于不同返回类型的函数重载。
\end{noteBox}

\subsection{函数递归}
\subsubsection*{Function Recursion}

PXL 不支持递归函数调用。例如：
\begin{lstlisting}[style=pxl]
noCanDo : function( ... ) returning integer {
  ...
  x = noCanDo(...);
  ...
}
\end{lstlisting}

% ============================================================
\section{程序结构}
\subsection*{Program Structure}

PXL 程序是声明、函数定义与语句的序列。程序自上而下执行。

PXL 对程序有效性设定下列前提条件：
\begin{itemize}
  \item 所有对象必须在使用前定义。
  \item 所有语句仅因副作用而求值。表达式必须是赋值的一部分。
\begin{lstlisting}[style=pxl]
10 + 5;     // illegal because the result (15) is not captured
\end{lstlisting}
\end{itemize}

\subsection{复合语句}
\subsubsection*{Compound Statement}

复合语句定义可声明局部变量的嵌套作用域。复合语句可用在允许单条语句的任何地方。

\begin{seeAlsoBox}
关于作用域，见「变量作用域」（原书第~197~页）。
\end{seeAlsoBox}

\begin{itemize}
  \item 复合语句允许局部变量隐藏其他作用域级别先前定义的变量。编译器报告警告。隐藏也称为遮蔽（shadowing）。
  \item 复合语句使用花括号（\cmd{\{\}}）标识。
\begin{lstlisting}[style=pxl]
{ ; }      // Smallest legal compound statement
\end{lstlisting}
  \item 复合语句遵循下列语法：
\begin{lstlisting}[style=pxl]
{
    s1;
    s2;
    ...
    sN;
}
\end{lstlisting}
\end{itemize}

\subsection{违例}
\subsubsection*{Violations}

某些 IC Validator 函数可产生违例数据（也称为错误输出），替代或补充层输出。这些违例存储在 SQLite 数据库（错误数据库，PYDB）中，可供 VUE 用于调试，并在运行结束时写入 \cmd{LAYOUT\_ERRORS} 文件。也有一些 PXL 函数以违例数据为输入，例如 \cmd{write\_gds()} 与 \cmd{violation\_empty()} 函数。

若希望直接访问 IC Validator 工具产生的错误数据，见附录~E「PYDB Perl API」，了解如何使用 IC Validator Perl 应用程序编程接口（API）。

\begin{seeAlsoBox}
附录~E，PYDB Perl API。
\end{seeAlsoBox}

\subsubsection{违例注释}
\subsubsection*{Violation Comments}

所有违例输出都与一条注释字符串关联。该注释作为标题出现在 \cmd{LAYOUT\_ERRORS} 文件以及 VUE 中。任意数量的函数可输出到同一违例注释。

当前及嵌套 PXL 作用域的违例注释使用 \cmd{@} 运算符指定。若 runset 中未指定注释，违例获得默认注释 \cmd{"Violation"}。

下例展示如何为层集合中的每一层编写唯一的违例注释。违例命名为 ``Metal1 spacing <= 0.1um''、``Metal2 spacing <= 0.1um''，依此类推。
\begin{lstlisting}[style=pxl]
for (i=0 to metal_layers.size()-1) {
   @ "Metal" + (i+1) + " spacing <= 0.1um";
   external1(metal_layers[i], distance<=0.1);
}
\end{lstlisting}

下例展示违例注释相对于作用域的行为。

\begin{noteBox}
通常，违例注释定义（\cmd{@} 运算符）不应与其所适用的函数分开。下例突出说明违例注释相对于作用域的行为。
\end{noteBox}

\begin{lstlisting}[style=pxl]
@ "First Error Message";
{ @ "Second Error Message";             // A new scope is created here
  local_var = abc(layer2);
  internal1(local_var, distance < 0.2); // Errors from internal1()
                                        // are stored in the database
                                        // under "Second Error Message"
} // The original scope is restored here

external1(layer1, distance < 0.1); // Errors from external1() are stored
                                   // in the database under "First Error
                                   // Message"
\end{lstlisting}

下例展示如何为层集合中的每一层编写唯一的违例注释。违例命名为 ``Metal1 spacing <= 0.1um''、``Metal2 spacing <= 0.1um''，依此类推。
\begin{lstlisting}[style=pxl]
for (i=0 to metal_layers.size()-1) {
   @ "Metal" + (i+1) + " spacing <= 0.1um";
   external1(metal_layers[i], distance<=0.1);
}
\end{lstlisting}

\cmd{text\_net()}、\cmd{texted\_with()} 与 \cmd{not\_texted\_with()} 函数可输出多种不同类型的违例。对这些函数，有办法为个别错误类型覆盖当前 runset 违例注释。本例中，text short 输出到违例注释 ``Rule 4.5.6: No Shorts''，而所有其他类型的 text 错误输出到 ``Rule 1.2.3: No text errors''。
\begin{lstlisting}[style=pxl]
@ "Rule 1.2.3: No text errors";

cdb1 = text_net(cdb1,
                { { M1, M1_text } },
                shorted_violation_comment = "Rule 4.5.6: No shorts");
\end{lstlisting}

IC Validator 工具使用诸如 \cmd{-svc} 等命令行选项处理违例注释。例如：
\begin{lstlisting}[style=pxl]
{@ "A";
   {@ "B";
      internal1(...);
   };
};
\end{lstlisting}

工具在选择时仅考虑较低层的违例注释。在此情况下，设置 \cmd{-svc B} 或 \cmd{-uvc A} 会执行 \cmd{internal1()} 函数；而设置 \cmd{-svc A} 或 \cmd{-uvc B} 则不会。

\subsubsection{产生违例的函数}
\subsubsection*{Violation-Producing Functions}

大多数 DRC 与数据生成函数是重载的；也就是说，它们有两个版本。返回多边形层的版本不产生错误输出。返回 \cmd{void} 的版本产生错误输出但不产生多边形层。下面是 \cmd{external1()} 函数两个版本的示例。
\begin{lstlisting}[style=pxl]
// This external1 outputs to a polygon layer, no errors are produced
t1 = external1(layer1, distance < 0.1);

// This external1 outputs error data to the comment "Ext1 Check"
@ "Ext1 Check";
external1(layer1, distance < 0.1);
\end{lstlisting}

也有一些 PXL 函数不返回 \cmd{void}，但也可产生错误输出。例如，\cmd{text\_net()} 函数返回 connect database，同时也为 text short、open 与 unused text 产生错误输出。

\subsubsection{违例变量}
\subsubsection*{Violation Variables}

内置 PXL 类型 \cmd{violation} 可用于保存错误数据的变量，以便稍后在 runset 中将其用作另一 PXL 函数的输入。例如，违例变量可用作 \cmd{write\_gds()} 等函数的输入以输出错误多边形，或用作 \cmd{violation\_empty()} 以判断违例是否包含任何错误。

如下例所示，某些 PXL 函数直接返回违例：
\begin{lstlisting}[style=pxl]
// These functions return violations for which there is not an associated
// PXL function
v1 = get_layout_grid_violation();
v2 = get_layout_drawn_violation();

// Write errors to GDSII
g = gds_library("err.gds");
write_gds(output_library = g,
          errors = { { v1, { 1 } },
                     { v2, { 2 } } });
\end{lstlisting}

\begin{noteBox}
错误输出可受 \cmd{error\_options()} 函数中 \cmd{error\_limit\_per\_check} 等实参限制。这些选项可影响实际存入错误数据库的数据量。被丢弃的错误对 \cmd{write\_gds()} 及其他函数不可用。
\end{noteBox}

\subsubsection{违例块}
\subsubsection*{Violation Blocks}

违例变量也可使用违例块赋值。违例块使用 \cmd{@=} 运算符指定。左侧为违例变量。\cmd{@=} 运算符后跟花括号括起的代码块。在该块内，所有产生错误输出的函数都与该违例变量关联。例如：
\begin{lstlisting}[style=pxl]
// v1 contains the error output from the external1() check
v1 @= {
   @ "Rule 1";
   l1 = a and b not c;
   external1(l1, distance < 2);
}

// v2 contains the error output from both internal1() checks and
// the external1() check
v2 @= {
   @ "Rule 2";
   internal1(l2, distance < 1);
   internal1(l3, distance < 1);

   @ "Rule 3";
   external1(l2, distance < 1);
}

tdb : connect_database; // Declared outside the violation block scope
                        // so it can be used after the violation block
// v3 contains all text errors found by text_net()
v3 @= {
   @ "Rule 4";
   tdb = text_net(cdb, {{ M1, M1_text }});
}

// These functions can be used to get only certain types of errors
// from a violation. The result is a subset of the input violation.
v4 = get_text_shorted(v3); // v4 contains only text shorts from v3
v5 = get_text_unused(v3); // v5 contains only unused text from v3

// Take some action if v4 is not empty
if (! violation_empty(v4)) {
   note("Found text shorts!");
}
\end{lstlisting}

\subsection{执行模型}
\subsubsection*{Execution Model}

PXL 的执行模型要求资源可用。若某程序中操作的前提条件已满足，则该操作随时可供执行。强烈偏向尽早执行诊断，意味着在程序初始处理期间，所有已就绪的诊断尽可能立即执行。因此，你可能看到 PXL 程序中语句乱序执行。最终结果不受此重排影响；但是，中间结果可能在不同时间出现或同时出现。这种乱序执行也称为并行执行。

同一程序在不同次执行时始终生成相同结果。不过，这些结果可能因资源可用性而以不同顺序出现。可能的运行时错误包括：
\begin{itemize}
  \item 试图使用尚未赋值的任何变量。
  \item 索引超出列表末尾，或以负索引索引。
  \item 以不存在元素的键值从 hash 读取。
\end{itemize}

\subsection{预处理器}
\subsubsection*{Preprocessor}

PXL 使用预处理器使程序的开发与执行更简便、更快速。标准预处理器模块集成在 PXL 编译器中。表~\ref{tab:pxl-preprocessor} 展示预处理器能力。

\begin{noteBox}
预处理器的行长度限制为 64K 字符。超出此限制通常可通过在空白或分隔符（如逗号 \cmd{,}）处插入回车来规避。
\end{noteBox}

\begin{longtable}{@{}>{\raggedright\arraybackslash}p{0.28\textwidth} p{0.64\textwidth}@{}}
\caption{标准预处理器模块}\label{tab:pxl-preprocessor}\\
\toprule
典型预处理能力 & 所用语法 \\
\midrule
\endfirsthead
\caption[]{标准预处理器模块（续）}\\
\toprule
典型预处理能力 & 所用语法 \\
\midrule
\endhead
\bottomrule
\endfoot
条件编译 & \cmd{\#if} \ldots\ \cmd{\#ifdef}、\cmd{\#ifndef}、\cmd{\#elif}、\cmd{\#else} \\
错误生成 & \cmd{\#error} \\
文件包含 & \cmd{\#include} 用于将附加文件内容插入 runset。文件名必须用双引号或尖括号括起。语法控制搜索路径。\newline
用尖括号包含 IC Validator 文件：\newline
\cmd{\#include <icv.rh>}\newline
按下列顺序搜索目录：\newline
1.~命令行上用 \cmd{-I} 选项指定的目录，按其指定顺序\newline
2.~\cmd{ICV\_INSTALL\_DIRECTORY/include}\newline
用双引号包含你自己的文件：\newline
\cmd{\#include "my\_include.rh"}\newline
按下列顺序搜索目录：\newline
1.~含有 \cmd{\#include} 语句的文件所在同一目录\newline
2.~命令行上用 \cmd{-I} 选项指定的目录，按其指定顺序\newline
3.~\cmd{ICV\_INSTALL\_DIRECTORY/include}\newline
指定绝对路径时，仅搜索该路径。\newline
可用 \cmd{\#pragma once} 防止文件被包含超过一次。任何含有 \cmd{\#pragma once} 的文件只会被包含一次，即使它出现在多个 \cmd{\#include} 语句中。\newline
例如，若 runset 目录结构如下：\newline
\cmd{Runset}\newline
\cmd{|-- rules.rs}\newline
\cmd{|-- Include}\newline
\cmd{\phantom{|-- }|-- assign.rs}\newline
\cmd{\phantom{|-- }|-- other\_assign.rs}\newline
要在 \cmd{rules.rs} 中包含 \cmd{other\_assign.rs}，使用\newline
\cmd{\#include "Include/other\_assign.rs"}\newline
要在 \cmd{assign.rs} 中包含 \cmd{other\_assign.rs}，使用\newline
\cmd{\#include "other\_assign.rs"} \\
位置跟踪 & \cmd{\#line} \\
标准宏支持 & \cmd{\#define} \ldots\ \cmd{\#undef}，包括记号构造（宏体中的 \cmd{\#\#} 二元运算符）与字符串化（宏体中的 \cmd{\#} 一元运算符）。也可用 \cmd{-D} 命令行选项定义变量。 \\
\end{longtable}

\begin{seeAlsoBox}
IC Validator 命令行选项（原书第~25~页）。\\
关于 IC Validator 主目录的更多信息，见 SolvNetPlus 上的 IC Validator Installation Notes。
\end{seeAlsoBox}

使用预处理器指令移除大块代码；也就是说，将该区段包在一对 \cmd{\#if 0} 与 \cmd{\#endif} 语句内。例如：
\begin{lstlisting}[style=pxl]
#if 0
// everything through the matching #endif is treated as a comment.
if (something == true)
{
   a line comment.
   do_something(useful = true);
}
#endif
\end{lstlisting}

下面是使用指令强制单次包含的示例。首次遇到此文件时 \cmd{MYFILE\_RH} 未定义。所含的 \cmd{\#define MYFILE\_RH} 将该变量设为已定义，并包含代码体。此后因 \cmd{MYFILE\_RH} 已存在且 \cmd{\#ifndef} 求值为 false，代码体不再被包含。
\begin{lstlisting}[style=pxl]
#ifndef MYFILE_RH
#define MYFILE_RH

// Body of code to include goes here

#endif        // end if the #ifndef MYFILE_RH
\end{lstlisting}

使用 \cmd{\#if} 语句时：
\begin{itemize}
  \item 使用标准比较运算符：\cmd{==}、\cmd{<=}、\cmd{>=}、\cmd{\&\&}、\cmd{||}
  \item 未定义的宏被解释为 0。例如，若 \cmd{SOMEMACRO} 未定义，则 \cmd{\#if SOMEMACRO == 1} 被翻译为 \cmd{\#if 0 == 1}。
  \item 仅比较整数值。不能比较字符串。
\end{itemize}

\subsubsection{预定义标识符}
\subsubsection*{Predefined Identifiers}

预处理期间支持表~\ref{tab:pxl-predef-id} 所示预定义标识符。

\begin{longtable}{@{}>{\raggedright\arraybackslash\ttfamily}p{0.36\textwidth} p{0.56\textwidth}@{}}
\caption{预定义标识符}\label{tab:pxl-predef-id}\\
\toprule
标识符 & 内容 \\
\midrule
\endfirsthead
\caption[]{预定义标识符（续）}\\
\toprule
标识符 & 内容 \\
\midrule
\endhead
\bottomrule
\endfoot
\_\_LINE\_\_ & 当前源文件中的行号，为整数 \\
\_\_FILE\_\_ & 当前源文件名，为字符串常量 \\
\_\_DATE\_\_ & 编译日期，为 11 字符字符串字面量，例如 ``Apr 25 2006'' \\
\_\_TIME\_\_ & 编译时间，为字符串字面量（``hh:mm:ss''） \\
\_\_PXL\_VERSION\_MAJOR\_\_ & 语言主版本号，为整数 \\
\_\_PXL\_VERSION\_MINOR\_\_ & 语言次版本号，为整数 \\
\_\_PXL\_VERSION\_PATCH\_\_ & 语言补丁级别，为整数 \\
\_\_PXL\_DATE\_\_ & 编译日期，为八位整数常量 yyyymmdd，例如 20060425 \\
\_\_PXL\_TIME\_\_ & 编译时间，为四位整数常量 hhmm，例如 0014 或 2300 \\
\end{longtable}

例如，\cmd{\_\_PXL\_DATE\_\_} 与 \cmd{\_\_PXL\_TIME\_\_} 可用于编写基于当前日期改变行为的代码。因为这些变量是预处理器变量，也可用于 \cmd{\#if} 与 \cmd{\#define} 语句。
\begin{lstlisting}[style=pxl]
/**
 * define a new token, expire_20061231, that is be replaced
 * by the token deprecated before 2006.12.31 but expands
 * to obsolete thereafter.
 */
#if __PXL_DATE__ < 20061231
#define expire_20061231 deprecated
#endif

#ifndef expire_20061231
#define expire_20061231 obsolete
#endif
\end{lstlisting}

\subsubsection{发布版本宏}
\subsubsection*{Release Version Macros}

已发布宏允许你将当前发布版本与其他发布版本比较。

例如，你可用年、月、service pack 与 hot fix 的整数值判断当前发布版本是否至少与指定发布版本一样新。
\begin{lstlisting}[style=pxl]
#if VERSION_GE(2009, 6, 1, 0)
. . .      // code that uses new features introduced in 2009.06.SP01
#else
. . .      // code that uses older method (pre 2006.06.SP01)
#endif
\end{lstlisting}

VERSION 宏定义见表~\ref{tab:pxl-version-macros}。

\begin{noteBox}
\cmd{year}、\cmd{month}、\cmd{service\_pack} 与 \cmd{patch} 为整数值。
\end{noteBox}

\begin{longtable}{@{}>{\raggedright\arraybackslash\ttfamily}p{0.42\textwidth} p{0.50\textwidth}@{}}
\caption{VERSION 宏}\label{tab:pxl-version-macros}\\
\toprule
宏 & 定义 \\
\midrule
\endfirsthead
\caption[]{VERSION 宏（续）}\\
\toprule
宏 & 定义 \\
\midrule
\endhead
\bottomrule
\endfoot
VERSION\_GE(year, month, service\_pack, patch) & 测试当前发布版本是否大于或等于指定发布版本。 \\
VERSION\_LE(year, month, service\_pack, patch) & 测试当前发布版本是否小于或等于指定发布版本。 \\
VERSION\_GT(year, month, service\_pack, patch) & 测试当前发布版本是否大于指定发布版本。 \\
VERSION\_LT(year, month, service\_pack, patch) & 测试当前发布版本是否小于指定发布版本。 \\
VERSION\_EQ(year, month, service\_pack, patch) & 测试当前发布版本是否等于指定发布版本。 \\
\end{longtable}

\subsubsection{日期比较宏}
\subsubsection*{Date Comparison Macros}

日期比较宏允许你根据当前日期改变 runset 行为。

例如，可用新宏按当前日期启用或禁用某些 PXL 代码。
\begin{lstlisting}[style=pxl]
if (DATE_EQ(2015, 6, 4)) {
    note("Today == 2015-06-04");
} else {
    note("Today != 2015-06-04");
}
\end{lstlisting}

DATE 宏定义见表~\ref{tab:pxl-date-macros}。

\begin{noteBox}
\cmd{year}、\cmd{month} 与 \cmd{day} 为整数值。指定无效或越界值（例如 day 值大于 31，或 month 值大于 12）会导致未定义行为。
\end{noteBox}

\begin{longtable}{@{}>{\raggedright\arraybackslash\ttfamily}p{0.36\textwidth} p{0.56\textwidth}@{}}
\caption{DATE 宏}\label{tab:pxl-date-macros}\\
\toprule
宏 & 定义 \\
\midrule
\endfirsthead
\caption[]{DATE 宏（续）}\\
\toprule
宏 & 定义 \\
\midrule
\endhead
\bottomrule
\endfoot
DATE\_EQ(year, month, day) & 测试当前日期是否等于由 year、month、day 整数指定的日期。 \\
DATE\_GE(year, month, day) & 测试当前日期是否大于或等于由 year、month、day 整数指定的日期。 \\
DATE\_GT(year, month, day) & 测试当前日期是否大于由 year、month、day 整数指定的日期。 \\
DATE\_LE(year, month, day) & 测试当前日期是否小于或等于由 year、month、day 整数指定的日期。 \\
DATE\_LT(year, month, day) & 测试当前日期是否小于由 year、month、day 整数指定的日期。 \\
\end{longtable}

\subsection{限定符}
\subsubsection*{Qualifiers}

限定符通过改变其所标记项的行为或允许的操作来影响该项。表~\ref{tab:pxl-qualifiers} 定义各限定符。

\begin{longtable}{@{}>{\raggedright\arraybackslash\ttfamily}p{0.22\textwidth} p{0.70\textwidth}@{}}
\caption{限定符}\label{tab:pxl-qualifiers}\\
\toprule
限定符 & 定义 \\
\midrule
\endfirsthead
\caption[]{限定符（续）}\\
\toprule
限定符 & 定义 \\
\midrule
\endhead
\bottomrule
\endfoot
const & 每次对该对象赋新值都会导致错误。初始化不导致此错误。仅在尝试直接赋值或通过带 \cmd{out} 或 \cmd{in\_out} 限定的函数参数传递时才会报错。 \\
deprecated & 每次定义或引用该对象都会导致警告。在函数原型中为函数参数赋初值，或在该参数所属函数体内使用该参数，不会生成此警告。 \\
in\_out & 函数参数在每次函数调用中要求可赋值项（lvalue）。不可赋值项或类型错误的项会导致错误。该项的值也会因函数调用而改变。关于 lvalue 的更多信息，见「赋值运算符」（原书第~199~页）。 \\
inline & 将函数标记为在使用 \cmd{-ndg} 命令行选项的 IC Validator 运行中需要展开。见 IC Validator 命令行选项（原书第~25~页）。 \\
literal & 函数参数要求函数实参为字面量表达式。也就是说，编译器必须能在解析时解析函数实参的值。 \\
obsolete & 每次定义或引用该项都会导致错误。在函数原型中为函数参数赋初值不会生成此错误。 \\
out & 函数参数在每次函数调用中要求可赋值项（lvalue）。不可赋值项或类型错误的项会导致错误。处理函数体时假定该参数无值。该项的值也会因函数调用而改变。关于 lvalue 的更多信息，见「赋值运算符」（原书第~199~页）。 \\
\end{longtable}

限定符只能用于限定相关项。限定符见表~\ref{tab:pxl-qualifier-items}。

\begin{longtable}{@{}>{\raggedright\arraybackslash\ttfamily}p{0.20\textwidth}
  >{\centering\arraybackslash}p{0.16\textwidth}
  >{\centering\arraybackslash}p{0.16\textwidth}
  >{\centering\arraybackslash}p{0.18\textwidth}
  >{\centering\arraybackslash}p{0.18\textwidth}@{}}
\caption{限定符适用项}\label{tab:pxl-qualifier-items}\\
\toprule
限定符 & 变量 & 函数 & 参数 & 结构成员 \\
\midrule
\endfirsthead
\caption[]{限定符适用项（续）}\\
\toprule
限定符 & 变量 & 函数 & 参数 & 结构成员 \\
\midrule
\endhead
\bottomrule
\endfoot
binary & & 是 & & \\
const & 是 & & 是 & \\
deprecated & 是 & 是 & 是 & 是 \\
in\_out & & & 是 & \\
inline & & 是 & & \\
literal & & & 是 & \\
obsolete & 是 & 是 & 是 & 是 \\
out & & & 是 & \\
\end{longtable}

% ============================================================
\section{模块化代码示例}
\subsection*{Example of Modularized Code}

\begin{itemize}
  \item \cmd{function.rh}
\begin{lstlisting}[style=pxl]
// Use pragma once to prevent including
// this file multiple times
#pragma once

#include <icv.rh>

// Function prototypes
fx1 : function (
   lyr : polygon_layer,
   w   : double
) returning out: polygon_layer;

fx2 : function (
   lyr : polygon_layer,
   d   : double
) returning out: polygon_layer;

#include <functions.rs>
\end{lstlisting}

  \item \cmd{function.rs}
\begin{lstlisting}[style=pxl]
// Function definitions
fx1 : function (
   lyr : polygon_layer,
   w   : double
   ) returning out: polygon_layer {
      // code body
      out = f( lyr, w);
   };

fx2 : function (
   lyr : polygon_layer,
   d   : double
   ) returning out: polygon_layer {
      // code body
      out = f( lyr, d);
   };
\end{lstlisting}

  \item \cmd{main.rs}
\begin{lstlisting}[style=pxl]
#include <icv.rh>
#include <function.rh>

// assign(), library(), etc.

y = fx1( a, 0.1);
z = fx2( b, 0.3);
\end{lstlisting}
\end{itemize}
